\section{Conclusion}

Under the frozen 11-dataset protocol and full six-architecture portfolio, the evaluated diagnostic rules do not demonstrate incremental value over always-graph in this benchmark. Combined with the declared MLP fallback incurs 7.46 points of regret versus 0.26 for always-graph; 75.9\% of its regret occurs on abstaining units. The fallback and dataset-exclusion sensitivities retain higher descriptive regret for Combined in the four specified configurations, while substantially changing the size of the difference. Post-hoc analysis reveals a concrete reversal: the unchanged rule has lower mean regret with GCN or GAT as the graph action, but higher regret with each other single architecture. Its favorable averages disappear after excluding both Cornell and Wisconsin. These observations show sensitivity within the evaluated candidate portfolios and dataset composition. Held-out threshold calibration reduces full-portfolio regret without matching always-graph, while GCN, GAT, and GPR-GNN have favorable calibrated averages; the findings do not establish general diagnostic failure.

The practical lesson is to evaluate the decision a diagnostic would make: specify its candidate portfolio and resource budget, compare its regret and coverage with relevant reference policies, and inspect baseline training and repeatability. Preprocessing changes the two-dataset test gaps, and validation-only diagnostics reveal training sensitivity and repeatability limits under the tested configurations. Our frozen artifacts and evaluation protocol support reconstruction of the recorded decision results; the bounded training audit does not replace a complete benchmark rerun or establish selection performance beyond this benchmark.
